\subsection{因果标量场 }
一般形式的场是：
\begin{equation}
    \Psi_{ab}(x) =(\frac{1}{2\pi})^\frac{3}{2}\sum_{\sigma}\int d^3p \, \left\{ u_{ab}(\boldsymbol{q},\sigma)b^\dagger(\boldsymbol{q},\sigma) e^{-ipx} + (-)^{2B}(-)^{j+\sigma}u_{ab}(\boldsymbol{q},\sigma) a(\boldsymbol{q},\sigma)e^{ipx} \right\}
\end{equation}
上述我们知道静止系系数 $u(0)$ 和 $v(0)$ 是正比于Clebsch-Gordan  系数 
\begin{equation}
    u_{ab}(0, \sigma) = \frac{1}{\sqrt{2m}} \langle A, a; B, b | j, \sigma \rangle
    \label{eq:u0_general}
\end{equation}
通过 Boost 获得有限动量系数 $u(\boldsymbol{q})$ 和 $v(\boldsymbol{q})$的一般公式是：
\begin{equation}
    u_{ab}(\boldsymbol{q}, \sigma) = \frac{1}{\sqrt{2q^0}} \sum_{\bar{a},\bar{b}} \left[e^{-\boldsymbol{\phi}\cdot\mathbf{J}^{(A)}}\right]_{a\bar{a}} \left[e^{+\boldsymbol{\phi}\cdot\mathbf{J}^{(B)}}\right]_{b\bar{b}} \langle A,\bar{a}; B,\bar{b} | j, \sigma \rangle
    \label{eq:u_finite_general}
\end{equation}
其中$\boldsymbol{\phi} = \hat{\boldsymbol{q}} \text{arctanh}(|\boldsymbol{q}|/q^0)$
\vspace{0.5cm}
\begin{center}
    \large\text{$(A, B) = (0, 0)$}
\end{center}
\vspace{0.5cm}
对于标量场，$(A, B) = (0, 0)$，粒子自旋 $j = 0$。此时：
\begin{itemize}
    \item 角动量生成元全为零：$\mathbf{J}^{(A)} = \mathbf{J}^{(B)} = 0$。
    \item 表示指标只有一个值：$a = 0$，$b = 0$。
    \item 自旋投影只有一个值：$\sigma = 0$。
    \item $j=0$根据自旋统计表明是波色子。
\end{itemize}
代入式 \eqref{eq:u0_general}，唯一的 C-G 系数为 $\langle 0, 0; 0, 0 | 0, 0 \rangle = 1$。因此静止系下的系数为：
\begin{equation}
    u(0, 0) = \frac{1}{\sqrt{2m}}
\end{equation}
\vspace{2ex}
对于标量场，$\mathbf{J}^{(A)} = 0$ 且 $\mathbf{J}^{(B)} = 0$。因此，指数演化矩阵完全退化为恒等映射：
\begin{equation}
    \exp(-\boldsymbol{\phi}\cdot\mathbf{0}) = 1, \quad \exp(+\boldsymbol{\phi}\cdot\mathbf{0}) = 1
\end{equation}
代入式 \eqref{eq:u_finite_general}，并利用 $\langle 0,0;0,0|0,0\rangle = 1$，有限动量系数极为简单：
\begin{equation}
    u(\boldsymbol{q}, 0) = \frac{1}{\sqrt{2q^0}}
\end{equation}
所以最终因果标量场是：
\begin{equation}
    \phi(x) =(\frac{1}{2\pi})^\frac{3}{2}\int \frac{d^3p }{\sqrt{2q^0}}\, \left\{ b^\dagger(\boldsymbol{q},0) e^{-ipx} + a(\boldsymbol{q},0)e^{ipx} \right\}
\end{equation}若粒子与反粒子不同（带电标量场，$a \neq b$），系统具有整体 $U(1)$ 规范对称性，对应守恒荷算符：
\begin{equation}
    Q = \int d^3p\left(a^\dagger(\boldsymbol{q})a(\boldsymbol{q}) - b^\dagger(\boldsymbol{q})b(\boldsymbol{q})\right)
\end{equation}
利用玻色对易关系验证：
\begin{equation}
    [Q, \phi(x)] = -\phi(x), \quad [Q, \phi^\dagger(x)] = +\phi^\dagger(x)
\end{equation}
$\phi$ 每作用一次使系统总电荷减少 $1$（湮没正粒子或产生反粒子），$\phi^\dagger$ 使总电荷增加 $1$。
为了后续计算传播子，需要考察系数函数的完备性。对于标量场，自旋 $j = 0$，只有一个极化态 $\sigma = 0$：
\begin{equation}
    \sum_{\sigma=0} u(\boldsymbol{q},\sigma)\, u^*(\boldsymbol{q},\sigma) = \frac{1}{2q^0}, \quad \sum_{\sigma=0} v(\boldsymbol{q},\sigma)\, v^*(\boldsymbol{q},\sigma) = \frac{1}{2q^0}
\end{equation}
